
\graphicspath{{Chapter4/Figs/}}

\setcounter{chapter}{7}
\chapter*{Capítulo 4} 
\addcontentsline{toc}{chapter}{Capítulo 4}
\setcounter{figure}{0}
\setcounter{section}{0}

{\LARGE Estudio bioinformático sobre ARN pequeños y su actividad en \textit{A. thaliana}.}

\section{Introducción}

En plantas existen distintos tipos de ARNs pequeños con distintos tamaño que varían de 20 a 25 nucleótidos, origen y mecanismo de acción.
Basado en las diferencias en su biogénesis y su modo de acción, los ARN pequeños de plantas con funciones regulatorias en la expresión génica han sido agrupados en cuatro clases diferentes
\begin{itemize}
	\item siARNs 	(small interfering RNAs) que es la clase más abundante. 
	\item miARNs (microRNAs). De 20 a 22 nt de longitud.
	\item ta-siARN (trans-acting siARNs)
	\item nat-siARNs y nat-miARNs (natural antisense siRNAs y miRNAs)
\end{itemize}

\section{Resultados y Discusión}

\subsection{Pipeline automatizado para el análisis genómico de ARN pqueños y su actividad in vivo en plantas}

\begin{table}[!htbp]
\centering
\small
\caption{Bibliotecas de ARN pequeños}
\label{my-label}
\begin{tabular}{clcccc}
\rowcolor[HTML]{ECF4FF} 
\hline
Código       & \multicolumn{1}{c}{Muestras de plantas}    & Secuencias totales & \begin{tabular}[c]{@{}c@{}}Lecturas que\\ mapean contra\\ el genoma\end{tabular} & \begin{tabular}[c]{@{}c@{}}Lecturas únicas\\ que mapean\\ contra el genoma\end{tabular} & t/r/sn/snoRNA \\ \hline
col\_8c\_1s  & Muestra 1 a 8°C   & 26594034           & 15235935                                                                         & 2181244                                                                                 & 7933811       \\ \hline
col\_8c\_2s  & Muestra 2 a 8°C   & 24503495           & 14064726                                                                         & 2073756                                                                                 & 7191654       \\ \hline
col\_22c\_1s & Muestra 1 a 22°C  & 13922338           & 6805121                                                                          & 952466                                                                                  & 5420744       \\ \hline
col\_22c\_2s & Muestra 2 a 22°C  & 23037685           & 12430935                                                                         & 1699845                                                                                 & 7575147       \\ \hline
col\_37c\_1s & Muestra 1 a 37°C  & 22273823           & 9109248                                                                          & 1388169                                                                                 & 10127406      \\ \hline
col\_37c\_2s & Muestra 2 a 37°C  & 24307772           & 12748308                                                                         & 1867528                                                                                 & 8511721       \\ \hline
\end{tabular}
\end{table}


Estudiar si se modifica la biogénesis y actividad de los small RNAs en distintas temperaturas.


\begin{figure}[htbp!] 
    \centering    
    \includegraphics[width=1\textwidth]{cantidad_lecturas.png}
    \caption[]{
    \textbf{}
   }
     \label{fig:secuencias_unicas}
\end{figure}


\begin{figure}[htbp!] 
    \centering    
    \includegraphics[width=1\textwidth]{cantidad_lecturas.png}
    \caption[]{
    \textbf{}
   }
     \label{fig:secuencias_unicas}
\end{figure}


\begin{figure}[htbp!] 
    \centering    
    \includegraphics[width=.5\textwidth]{sRNA_strategy.png}
    \caption[]{
    \textbf{}
   }
     \label{fig:sRNA_strategy}
\end{figure}





\subsection{Biogénesis y actividad de ARN pequeños de plantas en distintas temperaturas}
